\documentclass[11pt]{article}

\usepackage[margin=1in]{geometry}
\usepackage{amsmath,amssymb,amsthm}
\usepackage{graphicx}
\usepackage{booktabs}
\usepackage{subcaption}
\usepackage[numbers]{natbib}
\usepackage{xcolor}
\usepackage{float}
\usepackage{longtable}

\title{A Reduced WNT--RA--HOX Model for Colorectal Cancer Stemness}
\author{Amir, Nate, PK}
\date{June 2026}

\begin{document}

\maketitle

\begin{abstract}
This document records the dimensional and nondimensional WNT--RA--HOX model,
its biological interpretation, parameter realization, and planned scientific
machine-learning extensions. It is the repository copy of the supplied
Overleaf notes.
\end{abstract}

\section{Model Overview}

This model describes the regulatory interactions among the WNT, HOX, and
retinoic acid (RA) pathways in colorectal cancer. The seven state variables
are
\[
B(t)=\text{$\beta$-catenin},\qquad
P(t)=\text{APC},\qquad
H_5(t)=\text{HOXA5},
\]
\[
H_{13}(t)=\text{HOXA13},\qquad
M(t)=\text{MYC},\qquad
R(t)=\text{RA},\qquad
C(t)=\text{CYP26A1}.
\]

The model contains a pro-stemness loop,
\[
\beta\text{-catenin}\longrightarrow\text{MYC}
\longrightarrow\text{HOXA13}\longrightarrow\beta\text{-catenin},
\]
a differentiation loop,
\[
\text{RA}\longrightarrow\text{HOXA5}\longrightarrow\text{APC}
\dashv\beta\text{-catenin},
\]
and an RA homeostasis loop,
\[
\text{RA}\longrightarrow\text{CYP26A1}\dashv\text{RA}.
\]
Reduced APC functionality represents APC-mutant colorectal cancer cells.

\section{Dimensional Model}

\begin{align}
\frac{dB}{dt}
&=
a_BW
+v_{13}\frac{H_{13}^{n_H}}{K_{13}^{n_H}+H_{13}^{n_H}}
-d_BB-k_{PB}PB-v_5\frac{H_5B}{K_5+B},
\\
\frac{dP}{dt}
&=
a_P\frac{1+\rho_5H_5}{1+\rho_BB+\rho_{13}H_{13}}
-\Big(\gamma_{P0}+\gamma_{P1}(1-\Theta_P)\Big)P,
\\
\frac{dH_5}{dt}
&=
a_5+v_R\frac{R}{K_R+R}-d_5H_5
-v_M\frac{MH_5}{K_M+M},
\\
\frac{dH_{13}}{dt}
&=
a_{13}
+v_{B13}\frac{B^{n_B}}{K_{B13}^{n_B}+B^{n_B}}
+v_{M13}\frac{M^{n_M}}{K_{M13}^{n_M}+M^{n_M}}
-d_{13}H_{13},
\\
\frac{dM}{dt}
&=
a_M+v_{BM}\frac{B^{n_B}}{K_{BM}^{n_B}+B^{n_B}}-d_MM,
\\
\frac{dR}{dt}
&=
u_R(t)-d_RR-k_{CR}CR,
\\
\frac{dC}{dt}
&=
a_C+v_{RC}\frac{R}{K_{RC}+R}
+v_{BC}\frac{B^{n_B}}{K_{BC}^{n_B}+B^{n_B}}-d_CC.
\end{align}

The RA input is
\begin{align}
u_R(t)
&=
u_0
+A_d\left[1+\cos\left(\frac{2\pi t}{T_d}-\phi\right)\right]
\nonumber\\
&\quad+
\frac{D}{2}\left[
\tanh(q(t-t_1))-\tanh(q(t-t_2))
\right].
\end{align}

\section{Nondimensionalization}

Define
\[
b=\frac{B}{B_0},\quad
p=\frac{P}{P_0},\quad
h_5=\frac{H_5}{H_{50}},\quad
h_{13}=\frac{H_{13}}{H_{130}},
\]
\[
m=\frac{M}{M_0},\quad
r=\frac{R}{R_0},\quad
c=\frac{C}{C_0},\quad
\tau=d_Bt.
\]
Thus \(d/dt=d_B\,d/d\tau\).

\subsection{\(\beta\)-catenin equation}

Substitution and division by \(d_BB_0\) give
\[
\boxed{
\frac{db}{d\tau}
=
w+\eta_{13}\frac{h_{13}^{n_H}}
{\kappa_{13}^{n_H}+h_{13}^{n_H}}
-b-\lambda_Ppb-\lambda_5\frac{h_5b}{\kappa_5+b}.
}
\]
Here
\[
w=\frac{a_BW}{d_BB_0},\quad
\eta_{13}=\frac{v_{13}}{d_BB_0},\quad
\kappa_{13}=\frac{K_{13}}{H_{130}},
\]
\[
\lambda_P=\frac{k_{PB}P_0}{d_B},\quad
\lambda_5=\frac{v_5H_{50}}{d_BB_0},\quad
\kappa_5=\frac{K_5}{B_0}.
\]

\subsection{APC equation}

Choose
\[
P_0=\frac{a_P}{\gamma_{P0}},\qquad
\varepsilon_P=\frac{d_B}{\gamma_{P0}},
\qquad
\delta_{P1}=\frac{\gamma_{P1}}{\gamma_{P0}},
\qquad
\theta_P=\Theta_P.
\]
After absorbing the dimensional state scales into the \(\rho\) coefficients,
\[
\boxed{
\frac{dp}{d\tau}
=
\frac{1}{\varepsilon_P}
\left[
\frac{1+\rho_5h_5}{1+\rho_Bb+\rho_{13}h_{13}}
-\delta_P(\theta_P)p
\right].
}
\]
The APC mutation function is
\[
\boxed{
\delta_P(\theta_P)=1+\delta_{P1}(1-\theta_P).
}
\]
Normal APC function corresponds to \(\theta_P=1\), while
\(0<\theta_P<1\) represents partial APC loss.

\subsection{HOXA5 equation}

With \(H_{50}=a_5/d_5\), \(\varepsilon_5=d_B/d_5\), and the corresponding
dimensionless gains and thresholds,
\[
\boxed{
\frac{dh_5}{d\tau}
=
\frac{1}{\varepsilon_5}
\left[
a_5^*+\eta_R\frac{r}{\kappa_R+r}
-h_5-\eta_M\frac{mh_5}{\kappa_M+m}
\right].
}
\]

\subsection{HOXA13 equation}

\[
\boxed{
\frac{dh_{13}}{d\tau}
=
\frac{1}{\varepsilon_{13}}
\left[
a_{13}^*
+\eta_{B13}\frac{b^{n_B}}{\kappa_{B13}^{n_B}+b^{n_B}}
+\eta_{M13}\frac{m^{n_M}}{\kappa_{M13}^{n_M}+m^{n_M}}
-h_{13}
\right].
}
\]

\subsection{MYC equation}

The dimensional MYC equation is
\[
\boxed{
\frac{dM}{dt}
=
a_M+v_{BM}\frac{B^{n_B}}{K_{BM}^{n_B}+B^{n_B}}-d_MM.
}
\]
Using \(M=M_0m\), \(B=B_0b\), \(\varepsilon_M=d_B/d_M\),
\(\eta_{BM}=v_{BM}/(d_MM_0)\), and
\(\kappa_{BM}=K_{BM}/B_0\), the dimensionless equation is
\[
\boxed{
\frac{dm}{d\tau}
=
\frac{1}{\varepsilon_M}
\left[
a_M^*
+\eta_{BM}\frac{b^{n_B}}
{\kappa_{BM}^{n_B}+b^{n_B}}
-m
\right].
}
\]
If \(M_0=a_M/d_M\), then \(a_M^*=1\).

\subsection{Retinoic acid equation}

\[
\boxed{
\frac{dr}{d\tau}
=
\frac{1}{\varepsilon_R}
\left[
\mu_R(\tau)-r-\lambda_Ccr
\right],
}
\]
where
\[
\varepsilon_R=\frac{d_B}{d_R},\qquad
\lambda_C=\frac{k_{CR}C_0}{d_R},
\]
and
\[
\boxed{
\mu_R(\tau)
=
\mu_0
+A_R\left[
1+\cos\left(\frac{2\pi\tau}{T_R}-\phi\right)
\right]
+\frac{D_R}{2}
\left[
\tanh(q(\tau-\tau_1))-\tanh(q(\tau-\tau_2))
\right].
}
\]

\subsection{CYP26A1 equation}

\[
\boxed{
\frac{dc}{d\tau}
=
\frac{1}{\varepsilon_C}
\left[
a_C^*
+\eta_{RC}\frac{r}{\kappa_{RC}+r}
+\eta_{BC}\frac{b^{n_B}}{\kappa_{BC}^{n_B}+b^{n_B}}
-c
\right].
}
\]

\section{Final Nondimensional System}

\begin{align}
\frac{db}{d\tau}
&=
w+\eta_{13}\frac{h_{13}^{n_H}}
{\kappa_{13}^{n_H}+h_{13}^{n_H}}
-b-\lambda_Ppb-\lambda_5\frac{h_5b}{\kappa_5+b},
\\
\frac{dp}{d\tau}
&=
\frac{1}{\varepsilon_P}
\left[
\frac{1+\rho_5h_5}{1+\rho_Bb+\rho_{13}h_{13}}
-\delta_P(\theta_P)p
\right],
\\
\frac{dh_5}{d\tau}
&=
\frac{1}{\varepsilon_5}
\left[
a_5^*+\eta_R\frac{r}{\kappa_R+r}
-h_5-\eta_M\frac{mh_5}{\kappa_M+m}
\right],
\\
\frac{dh_{13}}{d\tau}
&=
\frac{1}{\varepsilon_{13}}
\left[
a_{13}^*
+\eta_{B13}\frac{b^{n_B}}{\kappa_{B13}^{n_B}+b^{n_B}}
+\eta_{M13}\frac{m^{n_M}}{\kappa_{M13}^{n_M}+m^{n_M}}
-h_{13}
\right],
\\
\frac{dm}{d\tau}
&=
\frac{1}{\varepsilon_M}
\left[
a_M^*+\eta_{BM}\frac{b^{n_B}}
{\kappa_{BM}^{n_B}+b^{n_B}}-m
\right],
\\
\frac{dr}{d\tau}
&=
\frac{1}{\varepsilon_R}
\left[\mu_R(\tau)-r-\lambda_Ccr\right],
\\
\frac{dc}{d\tau}
&=
\frac{1}{\varepsilon_C}
\left[
a_C^*+\eta_{RC}\frac{r}{\kappa_{RC}+r}
+\eta_{BC}\frac{b^{n_B}}{\kappa_{BC}^{n_B}+b^{n_B}}-c
\right].
\end{align}

\section{Stemness Measures}

The dynamic and equilibrium stemness indices are
\[
\boxed{
S(\tau)=
\frac{b(\tau)\left(1+\alpha_{13}h_{13}(\tau)\right)}
{\left(1+p(\tau)\right)\left(1+\alpha_5h_5(\tau)\right)}
}
\]
and
\[
\boxed{
S^*=
\frac{b^*\left(1+\alpha_{13}h_{13}^*\right)}
{\left(1+p^*\right)\left(1+\alpha_5h_5^*\right)}.
}
\]
Large values indicate proliferative, stem-like states; small values indicate
differentiated states.

\section{Expected Model Behavior}

\begin{enumerate}
\item \textbf{Normal tissue:} high APC, controlled \(\beta\)-catenin,
moderate MYC and HOXA13, and low stemness.
\item \textbf{APC-mutant tissue:} reduced APC functionality increases
\(\beta\)-catenin, MYC, and HOXA13 while suppressing differentiation.
\item \textbf{RA treatment:} RA induces HOXA5, enhances APC activity,
suppresses \(\beta\)-catenin, and lowers stemness.
\item \textbf{CYP-mediated resistance:} CYP26A1 accelerates RA degradation
and reduces the effectiveness of RA-based therapies.
\end{enumerate}

\section{Dimensional Interpretation}

Use the reference scales
\[
d_B=1~\mathrm{hr}^{-1},
\]
\[
B_0=0.20,\quad P_0=1.00,\quad H_{5,0}=0.80,\quad
H_{13,0}=0.30,\quad M_0=0.30,\quad R_0=0.60,\quad C_0=0.40.
\]
These reproduce the initial condition
\[
y_0=(0.20,1.00,0.80,0.30,0.30,0.60,0.40).
\]
Because \(d_B=1~\mathrm{hr}^{-1}\), one dimensionless time unit is one hour.
The notes interpret the treatment as an in-vitro exposure after an untreated
equilibration period. They contain both a \(40\)--\(88\) hour prose
interpretation and parameter tables using \(t_2=80\); this source preserves
that supplied record.

\begin{table}[H]
\centering
\caption{Reference scales.}
\begin{tabular}{lll}
\toprule
Quantity & Description & Value\\
\midrule
\(d_B\) & Reference degradation rate & \(1.00~\mathrm{hr}^{-1}\)\\
\(B_0\) & \(\beta\)-catenin scale & 0.20\\
\(P_0\) & APC scale & 1.00\\
\(H_{5,0}\) & HOXA5 scale & 0.80\\
\(H_{13,0}\) & HOXA13 scale & 0.30\\
\(M_0\) & MYC scale & 0.30\\
\(R_0\) & RA:RAR scale & 0.60\\
\(C_0\) & CYP26A1 scale & 0.40\\
\bottomrule
\end{tabular}
\end{table}

\begin{longtable}{lll}
\caption{Recovered dimensional parameter values.}\\
\toprule
Parameter & Interpretation / units & Value\\
\midrule
\endfirsthead
\toprule
Parameter & Interpretation / units & Value\\
\midrule
\endhead
\bottomrule
\endfoot
\(d_B\) & \(\beta\)-catenin degradation \((\mathrm{hr}^{-1})\) & 1.000\\
\(v_{13}\) & HOXA13 synthesis \((\mathrm{hr}^{-1})\) & 0.150\\
\(K_{13}\) & HOXA13 activation constant & 0.165\\
\(k_{PB}\) & APC inhibition coefficient \((\mathrm{hr}^{-1})\) & 1.600\\
\(v_5\) & HOXA5 synthesis \((\mathrm{hr}^{-1})\) & 0.325\\
\(K_5\) & \(\beta\)-catenin activation constant & 0.100\\
\(\gamma_{P0}\) & Basal APC degradation \((\mathrm{hr}^{-1})\) & 1.000\\
\(\gamma_{P1}\) & Enhanced APC degradation \((\mathrm{hr}^{-1})\) & 3.500\\
\(\rho_5\) & HOXA5 repression coefficient & 1.375\\
\(\rho_B\) & \(\beta\)-catenin repression coefficient & 5.500\\
\(\rho_{13}\) & HOXA13 activation coefficient & 4.333\\
\(d_5\) & HOXA5 degradation \((\mathrm{hr}^{-1})\) & 0.833\\
\(a_5\) & Basal HOXA5 production \((\mathrm{hr}^{-1})\) & 0.100\\
\(v_R\) & RA-induced HOXA5 synthesis \((\mathrm{hr}^{-1})\) & 1.667\\
\(K_R\) & RA activation constant & 0.240\\
\(v_M\) & MYC-induced HOXA5 synthesis \((\mathrm{hr}^{-1})\) & 1.667\\
\(K_M\) & MYC activation constant & 0.150\\
\(d_{13}\) & HOXA13 degradation \((\mathrm{hr}^{-1})\) & 1.000\\
\(a_{13}\) & Basal HOXA13 synthesis \((\mathrm{hr}^{-1})\) & 0.054\\
\(v_{B13}\) & \(\beta\)-catenin-induced HOXA13 synthesis & 0.285\\
\(K_{B13}\) & \(\beta\)-catenin activation constant & 0.100\\
\(v_{M13}\) & MYC-induced HOXA13 synthesis & 0.165\\
\(K_{M13}\) & MYC activation constant & 0.150\\
\(d_M\) & MYC degradation \((\mathrm{hr}^{-1})\) & 1.667\\
\(a_M\) & Basal MYC synthesis \((\mathrm{hr}^{-1})\) & 0.090\\
\(v_{BM}\) & \(\beta\)-catenin-induced MYC synthesis & 0.675\\
\(K_{BM}\) & \(\beta\)-catenin activation constant & 0.100\\
\(d_R\) & RA degradation \((\mathrm{hr}^{-1})\) & 2.500\\
\(u_0\) & Basal RA production \((\mathrm{hr}^{-1})\) & 0.525\\
\(A_d\) & Dietary RA input \((\mathrm{hr}^{-1})\) & 0.060\\
\(D\) & ATRA treatment amplitude \((\mathrm{hr}^{-1})\) & 2.250\\
\(T_d\) & Dietary RA period (hr) & 24\\
\(q\) & Treatment transition rate \((\mathrm{hr}^{-1})\) & 0.300\\
\(t_1,t_2\) & Treatment start/end (hr) & 40, 80\\
\(k_{CR}\) & CYP-mediated RA degradation \((\mathrm{hr}^{-1})\) & 5.313\\
\(d_C\) & CYP degradation \((\mathrm{hr}^{-1})\) & 1.250\\
\(a_C\) & Basal CYP synthesis \((\mathrm{hr}^{-1})\) & 0.040\\
\(v_{RC}\) & RA-induced CYP synthesis \((\mathrm{hr}^{-1})\) & 0.750\\
\(K_{RC}\) & RA activation constant & 0.300\\
\(v_{BC}\) & \(\beta\)-catenin-induced CYP synthesis & 0.750\\
\(K_{BC}\) & \(\beta\)-catenin activation constant & 0.100\\
\end{longtable}

\begin{longtable}{lll}
\caption{Dimensionless parameter values used in the simulations.}\\
\toprule
Parameter & Definition & Value\\
\midrule
\endfirsthead
\toprule
Parameter & Definition & Value\\
\midrule
\endhead
\bottomrule
\endfoot
\(\eta_{13}\) & \(v_{13}/(d_BB_0)\) & 0.75\\
\(\kappa_{13}\) & \(K_{13}/H_{13,0}\) & 0.55\\
\(\lambda_P\) & \(k_{PB}P_0/d_B\) & 1.60\\
\(\lambda_5\) & \(v_5H_{5,0}/(d_BB_0)\) & 1.30\\
\(\kappa_5\) & \(K_5/B_0\) & 0.50\\
\(\varepsilon_P\) & \(d_B/\gamma_{P0}\) & 1.00\\
\(\delta_{P1}\) & \(\gamma_{P1}/\gamma_{P0}\) & 3.50\\
\(\varepsilon_5\) & \(d_B/d_5\) & 1.20\\
\(\eta_R\) & \(v_R/(d_5H_{5,0})\) & 2.50\\
\(\kappa_R\) & \(K_R/R_0\) & 0.40\\
\(\eta_M\) & \(v_M/(d_5H_{5,0})\) & 2.50\\
\(\kappa_M\) & \(K_M/M_0\) & 0.50\\
\(\varepsilon_{13}\) & \(d_B/d_{13}\) & 1.00\\
\(\eta_{B13}\) & \(v_{B13}/(d_{13}H_{13,0})\) & 0.95\\
\(\kappa_{B13}\) & \(K_{B13}/B_0\) & 0.50\\
\(\eta_{M13}\) & \(v_{M13}/(d_{13}H_{13,0})\) & 0.55\\
\(\kappa_{M13}\) & \(K_{M13}/M_0\) & 0.50\\
\(\varepsilon_M\) & \(d_B/d_M\) & 0.60\\
\(\eta_{BM}\) & \(v_{BM}/(d_MM_0)\) & 1.35\\
\(\kappa_{BM}\) & \(K_{BM}/B_0\) & 0.50\\
\(\varepsilon_R\) & \(d_B/d_R\) & 0.40\\
\(\mu_0\) & \(u_0/(d_RR_0)\) & 0.35\\
\(A_R\) & \(A_d/(d_RR_0)\) & 0.04\\
\(D_R\) & \(D/(d_RR_0)\) & 1.50\\
\(T_R\) & \(d_BT_d\) & 24\\
\(q_\tau\) & \(q/d_B\) & 0.30\\
\(\tau_1,\tau_2\) & \(d_Bt_1,d_Bt_2\) & 40, 80\\
\(\lambda_C\) & \(k_{CR}C_0/d_R\) & 0.85\\
\(\varepsilon_C\) & \(d_B/d_C\) & 0.80\\
\(\eta_{RC}\) & \(v_{RC}/(d_CC_0)\) & 1.50\\
\(\kappa_{RC}\) & \(K_{RC}/R_0\) & 0.50\\
\(\eta_{BC}\) & \(v_{BC}/(d_CC_0)\) & 1.50\\
\(\kappa_{BC}\) & \(K_{BC}/B_0\) & 0.50\\
\(W\) & WNT stimulus & 0.80, 1.00, 1.50, 2.00\\
\(\Theta_P\) & APC reduction factor & 1.00, 0.75, 0.50, 0.25\\
\(n_B,n_M,n_H\) & Hill coefficients & 2\\
\(\alpha_5,\alpha_{13}\) & Stemness weights & 1\\
\end{longtable}

\section{Scientific Machine Learning Extensions}

The deterministic model has the form
\[
\frac{dy}{d\tau}=f(y,\theta),
\qquad
y=(B,P,H_5,H_{13},M,R,C).
\]

\subsection{Project 1: Bayesian PINNs for Parameter Uncertainty}

The inverse problem is
\[
p(\theta\mid D)\propto p(D\mid\theta)p(\theta).
\]
A biologically meaningful initial subset is
\[
\Theta=\{\eta_R,\eta_{BM},\eta_{B13},\lambda_C,\theta_P\}.
\]
The state network maps time to all seven states, and the loss combines data,
physics, and prior contributions:
\[
\mathcal L=\mathcal L_{\mathrm{data}}
+\mathcal L_{\mathrm{physics}}+\mathcal L_{\mathrm{prior}}.
\]
Outputs are posterior distributions, credible intervals, and uncertainty
bands.

\subsection{Project 2: Hybrid Mechanistic--Neural Models}

The general hybrid model is
\[
\frac{dy}{d\tau}
=f_{\mathrm{known}}(y,\theta)+f_{\mathrm{NN}}(y).
\]
The supplied notes proposed an additive MYC correction
\[
\frac{dM}{d\tau}
=
\frac{1}{\varepsilon_M}
\left[
a_M+\eta_{BM}\frac{B^{n_B}}{\kappa_{BM}^{n_B}+B^{n_B}}-M
\right]
+f_{\mathrm{NN}}(H_{13},R).
\]
The current repository additionally studies replacement hybrids in which a
closed-form regulatory relationship itself is learned. For the MYC term,
\[
\boxed{
\frac{dm}{d\tau}
=
\frac{1}{\varepsilon_M}
\left[a_M^*+f_{\mathrm{MYC}}(b)-m\right].
}
\]
For APC mutation,
\[
\boxed{
\delta_{P,\mathrm{NN}}(\theta_P)
=1+f_{\mathrm{APC}}(1-\theta_P),
}
\]
\[
\boxed{
\frac{dp}{d\tau}
=
\frac{1}{\varepsilon_P}
\left[
\frac{1+\rho_5h_5}{1+\rho_Bb+\rho_{13}h_{13}}
-\delta_{P,\mathrm{NN}}(\theta_P)p
\right].
}
\]
Small networks and function-magnitude regularization are used so the known
mechanistic structure remains dominant and interpretable.

\subsection{Project 3: Stochastic PINNs}

The deterministic model can be replaced by
\[
dy=f(y,\theta)\,d\tau+G(y)\,dW_\tau.
\]
A first model adds stochasticity only to RA and MYC:
\[
dR=f_R(y,\theta)\,d\tau+\sigma_R\,dW_R,
\qquad
dM=f_M(y,\theta)\,d\tau+\sigma_M\,dW_M.
\]
A state-dependent extension learns
\[
\sigma_R=\sigma_{\mathrm{NN}}(R,C).
\]
Euler--Maruyama gives
\[
y_{n+1}=y_n+f(y_n,\theta)\Delta\tau+G(y_n)\Delta W_n,
\qquad
\Delta W_n\sim\mathcal N(0,\Delta\tau).
\]
The outputs are trajectory ensembles, variance bands, and
\(\mathbb P(S>S_{\mathrm{crit}})\).

\section{Expected Contributions and Manuscript Outline}

The Bayesian project quantifies uncertainty, the hybrid project learns missing
or uncertain regulation while preserving biological structure, and the
stochastic project models cell-to-cell variability. Together they move the
model toward uncertainty-aware and biologically interpretable prediction.

The manuscript outline is:
\begin{enumerate}
\item Abstract (250 words maximum).
\item Introduction covering colorectal-cancer biology and PINNs.
\item Model schematic, assumptions, equations, and nondimensionalization.
\item Numerical SciPy results with biological interpretation of transients.
\item PINN algorithms, errors, convergence, sensitivity, identifiability, and
Bayesian uncertainty.
\item Discussion, limitations, biological significance, and future hybrid and
stochastic directions. Gene-level regulation is stochastic, and hybrid models
require real data to learn interactions absent from the reduced equations.
\end{enumerate}

\IfFileExists{references.bib}{
  \bibliographystyle{unsrt}
  \bibliography{references.bib}
}{}

\end{document}
